
\begin{exercise}[subtitle={Almost sure but not \(L^p\) convergence \Coffeecup}]
    \label{ex: as but not Lp}
    Let \(U\sim \uniform(0,1)\) and define
    \[
        X_n = 2^n \ind{[0, 2^{-n})}(U).
    \]
    Show that \(X_n \overset{\as}\to 0\) but \(X_n\) does not converge to \(0\)
    in \(L^1\) and hence in no other \(L^p\) space with \(p \ge 1\).
\end{exercise}

\begin{solution}
    First, we show that \(X_n \overset{\as}\to 0\). For any \(\epsilon > 0\)
    we have
    \[\begin{aligned}
        \sum_{n=0}^\infty \prob{\abs{X_n - 0} \ge \epsilon}
        &\le \sum_{n=0}^\infty \prob{X_n \ge \min\set{\epsilon, 1}}
        \\
        &= \sum_{n=0}^\infty\prob{U < 2^{-n}} = \sum_{n=0}^\infty 2^{-n} < \infty.
    \end{aligned}
    \]
    Consequently, by Corollary \ref{cor: characterization of almost sure convergence} we have \(X_n \overset{\as}\to 0\).
    To see that \(X_n\) does not converge to \(0\) in \(L^1\) we observe
    \[
        \E{\abs{X_n - 0}}
        = \E{X_n}
        = 2^n \prob{U < 2^{-n}} = 1
        \not\to 0.
        \qedhere
    \]
\end{solution}
